\documentclass[10pt,a4paper]{amsart}
\usepackage[margin=19mm]{geometry}
\usepackage{amsmath,amssymb,mathtools}
\usepackage[scr=rsfs,cal=euler]{mathalfa}
\usepackage{xfrac}
\usepackage[unicode,hidelinks,bookmarks=false]{hyperref}
\newcommand{\ie}{i.e.\ }
\newcommand{\cf}{cf.\ }
\newcommand{\resp}{respectively\ }
\newcommand{\Subsection}{\subsection}
\providecommand{\nobreakdash}{\nobreak\hbox{-}}
\setlength{\emergencystretch}{2em}
\multlinegap0pt
\usepackage{bm}
\usepackage{bbm}
\usepackage{dsfont}
\usepackage{xspace}
\usepackage{cancel}
\usepackage{mathtools}
\usepackage{amsthm}
\makeatletter
\@ifundefined{thm}{\newtheorem{thm}{Thm}}{}
\@ifundefined{pro}{\newtheorem{pro}[thm]{Proposition}}{}
\makeatother
\title[Idempotents and Powers of Ideals in Quandle Rings]{Idempotents and Powers of Ideals in Quandle Rings}
\author{Valeriy Bardakov, Mohamed Elhamdadi}
\date{Source: arXiv:2601.07057v3 (CC BY 4.0)}
\begin{document}
\maketitle

\noindent\emph{Source and licence.} Idempotents and Powers of Ideals in Quandle Rings. Version arXiv:2601.07057v3 (CC BY 4.0), distributed under the Creative Commons Attribution 4.0 International licence, as recorded in the arXiv licence metadata for this exact version and shown on the abstract page https://arxiv.org/abs/2601.07057v3. Full rights notice: licenses/arxiv-2601.07057-CC-BY-4.0.txt.\par\smallskip

\noindent\textit{Editorial scope.} Counted unit: the Pro whose source label is \texttt{pro4.8} in the article (source0.tex lines 505--525, source label \texttt{pro4.8}), delivered together with the article's own complete original proof of it (source0.tex lines 527--554, one \texttt{proof} environment of 28 lines). No mathematics is written, repaired or re-derived: statement and proof are the article's own text line for line. Required context retained uncounted: none. Every definition, display and label of those lines is retained unchanged. Omitted from this package and not redistributed: 1--504, 526--526, 555--1628. Nothing inside a retained excerpt is omitted or altered: every retained body is delivered exactly as the article's lines are written, so no omission receipt of any kind accompanies this package. Citations: the retained excerpts carry no citation command at all, so no bibliography entry of the article is transcribed and no reference is redistributed. Reference adaptation: none was needed and none was made; every reference inside the delivered unit resolves inside it, so no unresolved reference or citation is produced. Printed slips kept literal: no printed word, symbol, hypothesis or number of the counted statement or of its proof was altered. Layout and class: the article's own class is \texttt{amsart} with packages including amsfonts, amsbsy, amsmath, amssymb, fontenc, inputenc, lmodern, amsthm, amscd, mathrsfs, mathtools, tikz-cd, relsize, euler; the wrapper is the lane's pinned \texttt{amsart} editorial wrapper at 10pt on A4, which restates the theorem-like environments the retained text needs and adds the article's own macro definitions that the retained text needs (none), with extra wrapper packages (bm, bbm, dsfont, xspace, cancel, mathtools). The delivered numbering is the unit's own. Assets: no retained span contains a figure environment, a graphics-inclusion command, a PDF-page inclusion, a table or a TikZ picture; no image file is copied into this package, the package declares no asset dependency and no image file is redistributed with it. Licence: version arXiv:2601.07057v3 of this article is distributed under the Creative Commons Attribution 4.0 International licence, as recorded in the arXiv licence metadata for this exact version and shown on the abstract page https://arxiv.org/abs/2601.07057v3. A full rights notice is delivered at licenses/arxiv-2601.07057-CC-BY-4.0.txt. Characters: every retained excerpt is pure ASCII, so the delivered engine, XeLaTeX with its default Latin Modern text font, reports no missing character.
\begin{pro}\label{pro4.8}
For the dihedral quandle $R_3$, the powers of the augmentation ideal are given by:%we have in $\mathbb{Z}[R_3]$:

\[
\Delta(R_3)^{2k+1}
= \langle 3^k E_1,\; 3^k E_2 \rangle,
\quad k \ge 0,
\]
and
\[
\Delta(R_3)^{2\ell}
= \langle 3^{\ell-1}(E_1+E_2),\; 3^\ell E_2 \rangle,
\quad \ell \ge 1.
\]
 
 %$\Delta(R_3)^{2k + 1}$  has a basis: $\{ 3^k E_1, 3^k E_2 \}$, k = 0, 1, \ldots;


 %$\Delta(R_3)^{2l}$  has a basis: $\{ 3^{l-1} (E_1 + E_2), 3^l E_2 \}$, l = 1, 2, \ldots;

\end{pro}

\begin{proof}

Since $R_3$ is a commutative quandle, we see that $\Delta^2$ is generated by elements $(E_1)^2$, $E_1 E_2$, $(E_2)^2$. We find that
$$
(E_1)^2 = E_1 - 2 E_2,~~E_1 E_2 = - E_1 -  E_2,~~(E_2)^2 = E_2 - 2 E_1.
$$
It is easy to see that elements $E_1 + E_2$ and $3 E_2$ form a basis of $\Delta^2$.

Hence, we have a basis of induction.  By induction hypothesis we assume that for the odd powers,   $\Delta(R_3)^{2(k-1) + 1}$  has a basis: $\{ 3^{k-1} E_1, 3^{k-1} E_2 \}$. Then $\Delta(R_3)^{2k + 1} = \Delta(R_3)^{2(k-1) + 1} \cdot \Delta^2$ is generated by elements
$$
3^{k-1} E_1 \cdot (E_1 + E_2) = - 3^{k} E_2,~~3^{k-1} E_1 \cdot 3 E_2 = - 3^{k} (E_1 + E_2),
$$
$$
3^{k-1} E_2 \cdot (E_1 + E_2) =  - 3^k E_1,~~
3^{k-1} E_2 \cdot 3 E_2 =  3^k (-2 E_1 + E_2),
$$
and it is easy to check that the set  $\{ 3^{k} E_1, 3^{k} E_2 \}$ form a basis of  $\Delta(R_3)^{2k + 1}$.

 By  induction hypothesis we assume for the even powers,   $\Delta(R_3)^{2(l-1)}$, $l > 1$,  has a basis: $\{ 3^{l-2} (E_1 + E_2), 3^{l-1} E_2 \}$. Then $\Delta(R_3)^{2l} = \Delta(R_3)^{2(l-1)}\cdot \Delta^2$ is generated by the elements
$$
3^{l-2} (E_1 + E_2) \cdot (E_1 + E_2) = - 3^{l-1} (E_1 + E_2),~~3^{l-2} (E_1 + E_2) \cdot 3 E_2 = - 3^{l} E_1,
$$
$$
3^{l-1} E_2  \cdot (E_1 + E_2) =  - 3^l E_1,~~3^{l-1} E_2  \cdot 3 E_2 =  3^l (-2 E_1 + E_2),
$$
and it is easy to check that the set  $\{ 3^{l-1}(E_1 + E_2), 3^{l} E_2 \}$ form a basis of  $\Delta(R_3)^{2l}$.

\end{proof}

\end{document}
